\chapter*{Abstract}
\addcontentsline{toc}{chapter}{\protect\numberline{}Abstract}
Single-cell RNA-seq and single-cell ATAC-seq enable cell type atlases across organisms, tissues, and disease. To scale these efforts, lower cost, and support custom multiome assays, we built HyDrop and PUMATAC and systematically quantified method biases on a shared reference.

HyDrop is a flexible, open-source droplet platform for scATAC-seq and scRNA-seq based on dissolvable polyacrylamide hydrogel beads. Beads carry cycled split-pool barcodes and disulfide-cleavable BAC crosslinks. We redesigned the barcode layout for paired-end capture, removed custom priming, dissolved beads in-droplet to overcome diffusion and emulsion instability, tuned primer load to prevent cross-barcoding, and added Pitstop-2 and guanidinium to increase yield. HyDrop-ATAC generated \num{7996} high-quality single-cell profiles from flash-frozen mouse cortex. HyDrop-RNA produced \num{9508} transcriptomes from the same tissue and matched reference scRNA-seq. We further captured a small population of FACS-sorted \textit{Drosophila} neurons, demonstrating applicability to low-input samples. Across mouse cortex, HyDrop-RNA reached a median \num{3389} UMIs and \num{1658} genes per cell, and HyDrop-ATAC recovered a median \num{4148} unique fragments per cell with high TSS enrichment. Chemistry, microfluidic chips, and protocols are open-source and low-cost.

PUMATAC is a universal Nextflow preprocessing framework that unifies trimming, whitelist-aware barcode correction, alignment, and fragmentization across heterogeneous scATAC-seq formats. The pipeline integrates pycisTopic/LSI counting, Otsu-based cell calling, Scrublet doublet removal, Freemuxlet demultiplexing, MACS2 peak calling, and cisTarget motif analysis with containerization and restartable execution, and ships with best-practices code.

We benchmarked eight scATAC-seq methods on human PBMCs: 10x v1, v1.1, v2, 10x Multiome, Bio-Rad ddSEQ, mtscATAC-seq, s3-ATAC, and HyDrop. Across \num{47} libraries and \num{169227} high-quality cells, sequencing efficiency, defined as reads retained in the final cell-by-peak matrix, ranged from \textasciitilde\num{4}% for s3-ATAC and early HyDrop to \textasciitilde\num{28}% for 10x v2. Per-cell sensitivity scaled with efficiency at fixed depth and impacted cell-type annotation, genotype demultiplexing, peak calling, differential accessibility, and transcription factor motif enrichment. Methods also differed in specificity. Promoter and TSS proximity as well as FRIP varied markedly, with ddSEQ more TSS-proximal than 10x, yielding technique-dependent DARs and imperfect cross-technology integration even after batch correction. These results support a two-phase design in practice. Generate consensus peak sets with high-specificity kits, then scale cell numbers with cost-efficient methods at shallower depth.

Together, this work contributes HyDrop, a transparent high-performance droplet assay for scATAC-seq and scRNA-seq, PUMATAC, a reproducible method-agnostic preprocessing stack, and a public PBMC resource of \num{169000}+ scATAC-seq profiles with quantitative guidance on sensitivity, specificity, and depth-versus-cells trade-offs. All designs, code, and datasets are released for reuse.